\section{Related Work}
\label{sec:related_work}

We review the related work from the following five aspects:

\textbf{Multi-Key Homomorphic Encryption in FL:} HE provides a mathematically rigorous foundation for privacy-preserving FL. Single-key CKKS schemes, such as FedSHE~\cite{pan2024fedshe} and FedPHE~\cite{li2025fedphe}, support approximate arithmetic and ciphertext packing, but rely on a globally shared secret key, leaving local updates vulnerable to honest-but-curious clients. MKHE mitigates this risk by encrypting updates under independent client keys~\cite{chen2019efficient}, \cite{damgaard2012multiparty}, \cite{wu2025secure}. Notably, ESAFL~\cite{wu2025esafl} encodes packed message polynomials into a single ring element and achieves non-interactive \emph{one-step decryption} without uploading decryption shares, but it strictly requires full client participation. Conversely, MASER~\cite{al2025maser} aligns pruned parameter slices via majority voting under multi-key CKKS yet still requires interactive partial-decryption uploads, while FedOut~\cite{yuan2026fedout} achieves dropout resilience via two non-colluding online servers. To handle client dropouts, SecFed~\cite{cai2023secfed} delegates multi-key BFV key management and offline dropout compensation to a hardware TEE, whereas SAMK~\cite{li2025samk} recovers from dropouts by encrypting key shares under all participants' Paillier public keys at an $\mathcal{O}(N^2)$ setup cost. Overall, existing MKHE schemes either mandate full participation, rely on auxiliary hardware or twin servers, or incur interactive decryption-share and $\mathcal{O}(N^2)$ setup overhead, failing to simultaneously support client subsampling, upload-phase dropouts, and non-interactive one-step decryption.

\textbf{Threshold Homomorphic Encryption in FL:} Threshold HE employs Shamir secret sharing to distribute secret keys so any qualified subset can restore missing key contributions. Mouchet \emph{et al.}~\cite{mouchet2023threshold} instantiated RLWE threshold access structures over $\mathcal{R}_q$ using constant integer evaluation points, while ATASSES~\cite{chang2025atasses} handled non-participating dropout devices via approximate secret sharing with $\mathcal{O}(N^2)$ pairwise encrypted shares and CDKS-LSS~\cite{wang2026cdks} studied exceptional sequences over cyclotomic rings to guarantee invertible Lagrange interpolation. More recently, Colin de Verdi\`ere \emph{et al.}~\cite{colin2026tfhesd} formalized synchronous threshold FHE and proved that unmasked Lagrange weighting leaks secret keys across varying participant sets; Kim \emph{et al.} ($\Pi_{\mathrm{StThFHE}}$)~\cite{kim2026tfhead} mitigated this leakage using ring-based Shamir zero-sharing in the asynchronous setting. However, these threshold schemes all perform key restoration during interactive online decryption, exposing the plaintext aggregate to the server and incurring extra communication rounds. Embedding masked threshold key-contribution restoration directly into upstream encryption to preserve client-side one-step decryption remains unexplored.


\textbf{Gradient Compression:} To alleviate the communication bottleneck in FL, various gradient compression techniques have been proposed, including quantization~\cite{alistarh2017qsgd}, \cite{bao2025efskip}, low-rank decomposition~\cite{richtarik2024error}, \cite{dorfman2023docofl}, and sparsification~\cite{haddadpour2021federated}, \cite{kumar2023noniid}, \cite{kim2024spafl}. Among these, Top-$k$ sparsification with residual error feedback~\cite{stich2018sparsified} has proven highly effective. For encrypted FL, DSVFL~\cite{li2025enhancing} introduced FTop-$k$ with an accumulable residual buffer under Paillier encryption. However, element-wise sparsification before CKKS encryption destroys SIMD slot alignment across clients and leaks local gradient structures through sparse index bitmaps~\cite{xu2025harnessing}.


\textbf{Downstream Prefetching:} Under intermittent client participation, synchronizing stale local models often offsets the bandwidth savings of upstream compression. To mitigate this delay, FedFetch~\cite{yan2025fedfetch} presamples future candidates with over-commitment and schedules background downloads of compressed updates during idle windows. However, FedFetch operates exclusively on plaintext updates, addresses only pre-round unavailability rather than upload-phase dropouts after active-set confirmation, and cannot homomorphically accumulate or decrypt multi-round ciphertexts in a single step.

\textbf{Client Selection:} Client selection governs how the server samples the active subset per round, directly impacting convergence speed and wall-clock training time. While uniform Random selection~\cite{mcmahan2017communication} is vulnerable to stragglers, heterogeneity-aware schemes incorporate targeted selection criteria, including Power-of-Choice (PoC)~\cite{cho2022towards}, tier-based selection (TiFL)~\cite{chai2020tifl}, clustered selection (HACCS)~\cite{wolfrath2022haccs}, and utility-guided selection (Oort)~\cite{lai2021oort}. GlueFL~\cite{he2023gluefl} further introduces Sticky sampling to prioritize recently active clients with fresh models, and FedPHE~\cite{li2025fedphe} combines CKKS with locality-sensitive hashing (LSH)-based Sketch-Clustering. Yet, biased selection can skew global representation under non-IID data, and existing encrypted selection methods still force newly sampled clients to synchronously download full encrypted models upon joining.
